\section{Preliminaries}
\label{sec:preliminaries}

\subsection{Vision Language Action Model}

A deployed VLA policy operates in closed loop with a physical
environment~\citep{rt2,openvla,pi0}.
At each time step $t$, the policy receives the current observation
$o_t$, updates its internal hidden state $h_t$, and outputs an action
$a_t$ to be executed by the robot.
We write this runtime process as
\[
  h_t = F_\pi(h_{t-1}, o_t), \qquad a_t = G_\pi(h_t),
\]
where $\pi$ denotes the deployed policy and $h_t$ denotes the runtime
internal representation exposed by the model.
This hidden state is the primary runtime signal used by the fast layer
in the rest of the paper.

\subsection{General Pipeline of VLA Backdoor Attacks}

VLA backdoor attacks follow a common deployment pipeline even though
their concrete trigger types differ~\citep{badnets,goba,tabvla,contextbackdoor,statebackdoor}.
An attacker first poisons training data, injects a malicious
fine-tuning stage, or directly prepares compromised model weights, and
thereby obtains a backdoored policy $\pi'$.
After deployment, the compromised policy behaves similarly to the clean
policy on ordinary inputs, but when a trigger condition is satisfied it
can shift execution toward attacker-desired outcomes.
For deployment-time defense, the observable object is therefore the
resulting execution trajectory rather than the hidden poisoning process
that produced $\pi'$.

We formalize the standard backdoor semantics with a trigger indicator
$c(\tau)\in\{0,1\}$ on trajectory $\tau$.
When the trigger is absent, the compromised policy should preserve
benign task behavior:
\[
  P_{\pi'}(\tau \in \mathcal{S}_{\mathrm{task}} \mid c(\tau)=0)
  \approx
  P_{\pi}(\tau \in \mathcal{S}_{\mathrm{task}}).
\]
When the trigger is present, the compromised policy should induce
attacker-desired harmful behavior:
\[
  P_{\pi'}(\tau \in \mathcal{S}_{\mathrm{adv}} \mid c(\tau)=1)
  \ \text{is high},
\]
where $\mathcal{S}_{\mathrm{task}}$ denotes task-consistent execution
trajectories and $\mathcal{S}_{\mathrm{adv}}$ denotes attacker-desired
harmful trajectories.

\subsection{Threat Model}

We consider an attacker whose goal is to preserve clean task behavior
while causing harmful behavior under trigger activation.
The attacker may obtain a backdoored VLA through data poisoning,
malicious fine-tuning, or compromised model weights, but does not need
to retain control over the robot controller after deployment.
This threat model covers visual, language-conditioned, and
state-triggered VLA backdoors studied in recent
work~\citep{goba,tabvla,contextbackdoor,statebackdoor}.

The defender operates only at deployment time.
The defense can access runtime-visible signals such as observations,
hidden states, action streams, and their trajectory prefixes, but it cannot
access the poisoned training data, optimization history, or any future
information beyond the current execution prefix.
The problem studied here is therefore inference-time runtime defense
rather than training-time backdoor removal.

\subsection{Problem Formulation}

Our main task is backdoor-triggered dangerous-action detection.
We collect rollout trajectories generated by the deployed policy and
write each rollout as
\[
  \tau = (o_{1:T}, h_{1:T}, a_{1:T}),
\]
where $T$ is the rollout horizon.
Because the monitor runs online, it may only use the rollout prefix
$\tau_{1:t}$ available up to time step $t$.
For the primary task, each rollout is assigned one family-level label
from
\[
  \mathcal{Y}_{\mathrm{main}} =
  \{\mathtt{task\_success}, \mathtt{attack\_GoBA},
  \mathtt{attack\_Drop}, \mathtt{attack\_State}\}.
\]
We reserve
\[
  \mathcal{Y}_{\mathrm{aux}} = \{\mathtt{task\_fail}\}
\]
for auxiliary out-of-distribution failure experiments.
We then define the main binary target
\[
  z_{\mathrm{main}}(\tau) =
  \begin{cases}
    0, & \tau \in \mathtt{task\_success},\\
    1, & \tau \in \mathtt{attack\_GoBA}\ \cup\
    \mathtt{attack\_Drop}\ \cup\
    \mathtt{attack\_State}.
  \end{cases}
\]
This primary task asks whether the current rollout should be regarded
as unsafe to continue because it is heading toward a backdoor-triggered
dangerous action.

The auxiliary \mathtt{task\_fail} set is not part of the main target;
it is used only to probe whether the learned signals transfer to
non-adversarial failures.

The runtime detector is an online detector that maps the current
rollout prefix to a scalar risk score
\[
  q_t = D(\tau_{1:t}) \in [0,1].
\]
We define the first alarm time as
\[
  t^\star(\tau) = \min \{ t : q_t \ge \theta \},
\]
with $t^\star(\tau)=+\infty$ if no alarm is raised.
The deployment objective is to detect backdoor-induced dangerous
rollouts as early as possible while controlling false alarms on
successful rollouts:
\[
  \min_D \ \mathbb{E}\!\left[t^\star(\tau)\mid z_{\mathrm{main}}(\tau)=1\right]
  \
  \mathrm{s.t.}
  \
  \Pr\!\left(t^\star(\tau)<\infty \mid z_{\mathrm{main}}(\tau)=0\right)\le \alpha.
\]
Here $\alpha$ controls the tolerated false-alarm rate on successful
rollouts, and the next section describes the two-layer system used to
optimize this deployment objective.
